% CLAIM (reframed 2026-08-21): what the instrument must guarantee for the
% measurements to mean anything. Same four design paragraphs; their job is now
% validity, not product description. Evidence: src/, docs/02; E10.
\section{The instrument}
\label{sec:design}

Every measurement in this paper is a comparison, and a comparison is only as
good as the guarantee that both sides compute the same thing. This section
describes the instrument (shardes, built on JAX and its
compiler-driven sharding \citealp{jax2018,gspmd2021}) in those terms: not what the library
offers a user, but what it must guarantee for Sections~\ref{sec:results}
and~\ref{sec:endtoend} to mean anything. It exposes the classic
ask/evaluate/tell loop: \emph{ask} for $N$ perturbed copies of the model (the
population), \emph{evaluate} each copy's fitness, \emph{tell} the results
back to produce an update. Four properties do the work.

\paragraph{One mesh for every strategy.} Each of
$D$ devices holds the full weights and $N/D$ of the perturbed copies (the
population is sharded, the distribution state replicated). For isotropic ES
the alternative of splitting the distribution state saves nothing worth its
synchronization: isotropic or diagonal state is no larger than the
already-replicated model, while sharding it on the population axis would make
every device gather the same state again before perturbing every parameter
leaf.

\paragraph{Device count cannot change what is sampled or contracted.}
The seed for copy $i$ derives from $i$ alone (member-indexed seeds), so the
population is device-count invariant by construction, and given a fixed set
of fitnesses the update contracted on 1 device and on 8 agrees to
floating-point tolerance. This is a testable property and it is the most
important test in the repository. The guarantee is deliberately narrower
than trajectory equality: a full training trajectory is reproducible only
for the same compiled program, because compilation differences flip
near-ties in generation (Section~\ref{sec:endtoend} states the measured
boundary).

\paragraph{The cost comparison is fair to both bets.} A rank-$r$ perturbation is
carried as its two thin factors all the way through evaluation: the model's
dense layers apply $x(W + \tfrac{\sigma}{\sqrt r}AB^\top)^\top$ as
$xW^\top + ((xB)A^\top)\tfrac{\sigma}{\sqrt r}$, so the full $m\times n$ noise
matrix never exists in memory. The library treats a profile showing an
$(N, m, n)$ allocation under the low-rank path as a bug, and
Section~\ref{sec:results} shows what honoring that invariant buys: the dense
baseline runs out of memory over half of the measured grid while the
structured strategies keep running. The honest cost of this generality is
an API constraint: a model must route its parameterized matrix multiplies
and embedding lookups through the library's seams (\texttt{shardes.nn}),
because a structured weight substituted into the parameter tree raises on
direct arithmetic. Section~\ref{sec:limitations} carries the portability
consequence.

\paragraph{One deliberate synchronization point.} Fitness shaping needs every
candidate's score (a global sort for centered ranks), so \emph{tell} gathers
the $N$ scalars to every device before shaping: one line of code, one
barrier. We priced it directly (Section~\ref{sec:results}): the gather costs
2.4--6\,$\mu$s at every measured population and device count; the sort that
follows costs 12\,ms at $N{=}2^{18}$ and is identical on 1 and 8 devices,
i.e.\ it is redundant local compute, not communication. Making the gather
conditional on the shaping would save microseconds and cost the shaping
abstraction; we document the trade and decline it.
